% Generated draft: requires mathematical and source review.
\documentclass[11pt]{article}
\usepackage[margin=25mm]{geometry}
\usepackage{fontspec}
\setmainfont{Latin Modern Roman}
\usepackage{amsmath,amssymb,mathtools,unicode-math}
\setmathfont{Latin Modern Math}
\usepackage{enumitem,needspace,xurl,hyperref}
\setlength{\parindent}{0pt}
\setlength{\parskip}{5pt}
\setlength{\emergencystretch}{3em}
\newcommand{\N}{\mathbb{N}}
\newcommand{\Z}{\mathbb{Z}}
\newcommand{\Q}{\mathbb{Q}}
\newcommand{\R}{\mathbb{R}}
\newcommand{\C}{\mathbb{C}}
\newcommand{\F}{\mathbb{F}}
\newcommand{\A}{\mathbb{A}}
\newcommand{\PP}{\mathbb{P}}
\newcommand{\E}{\mathbb{E}}
\newcommand{\eps}{\varepsilon}
\newcommand{\dd}{\,\mathrm{d}}
\begin{document}
\section*{UnsolvedMath ID: 10600014}
\textbf{Source problem code:} catalogue entry
\subsection{Definitions and mathematical statement}
Fix a commutative coefficient ring \(R\) and loop parameter \(\delta\in R\). The Brauer algebra \(B_n(\delta;R)\) has an \(R\)-basis of all pairings of \(n\) upper and \(n\) lower endpoints. Multiplication concatenates diagrams and replaces each closed middle loop by a factor \(\delta\). The planar pairing diagrams form the Temperley--Lieb subalgebra. The source identifies its virtual Temperley--Lieb algebra, obtained by adjoining virtual crossings, with this algebra of all connections; this fixes the intended diagrammatic model rather than every possible algebra bearing that name.

Where the required denominators are invertible, the classical Jones--Wenzl element \(f_n\) is a normalized idempotent in the Temperley--Lieb algebra, annihilating the elementary turnback generators on both sides. Placing it in \(B_n\) preserves its idempotent identity under inclusion. The source's first question asks for the structure of these projectors in the Brauer setting, including whatever extensions or projector decomposition the intended virtual theory requires. It does not state new annihilation conditions for virtual crossings or specify that the full collection of Brauer primitive idempotents is the desired answer.

The second question asks whether a useful algebraic generalization of classical recoupling theory can be formulated in this setting. Such a formulation must specify its projectors, tensor/composition operations, decomposition and change-of-coupling data and the identities required for its intended application. The source does not prescribe these data or a mathematical criterion for usefulness. The ring, characteristic, \(\delta\), semisimplicity regime, projector normalization and generic versus specialized/root-of-unity treatment are missing. These choices affect the existence and behavior of projectors, so no unrestricted all-parameters assertion is added.
\subsection{Short English statement}
Study Jones--Wenzl projectors inside the Brauer algebra of all pairings and develop a corresponding virtual recoupling theory. The parameter regime and the exact projector or recoupling structure sought are unspecified.
\subsection{Sources}
\begin{itemize}
\item Local UnsolvedMath registry, permanent ID \(10600014\), exact source code \texttt{catalogue entry}.
\item \emph{Unsolved Problems in Virtual Knot Theory and Combinatorial Knot Theory}, revised survey, problem-list item 14: \url{https://arxiv.org/abs/1409.2823}. Registry access 2026-07-29, source-tex extraction. The embedded algebra citation is elided.
\item Supplied research leads: \emph{Virtual Extension of Temperley--Lieb Algebra} (2006; ``arXiv:0610.xxxx''); \emph{Projectors in the Virtual Temperley-Lieb Algebra} (2021; ``arXiv:2103.xxxx''); \emph{Anyonic Topological Quantum Computation and the Virtual Braid Group} (2009; ``arXiv:0909.xxxx''). Authors and valid full identifiers are not supplied; these placeholder identifiers are not usable citations.
\end{itemize}
\end{document}
